\section{Reto final integrador}

\subsection{Telemedicina 2.0 como caso transversal}

El reto integra problema, usuario, tarea, evidencia, característica, métrica, resultado, riesgo y recomendación en una sola cadena de trazabilidad. Se eligió la teleconsulta porque combina experiencia del paciente, operación clínica, infraestructura, contexto doméstico y continuidad asistencial.

El problema observado es la interrupción de sesiones. El usuario principal es el paciente, aunque el médico comparte el objetivo. La tarea consiste en completar una consulta remota con audio, video e indicaciones finales. El contexto incluye un dispositivo personal y una red fuera del control del hospital. La evidencia es el evento \texttt{teleconsultation}, que conserva inicio, sede asociada, rol, escenario, resultado y duración.

\begin{table}[H]
\centering
\caption{Trazabilidad del reto integrador}
\begin{tabularx}{\textwidth}{p{3.1cm}Y}
\toprule
Elemento & Definición aplicada \\
\midrule
Problemática & Sesiones interrumpidas antes de finalizar la atención. \\
Usuarios & Paciente y médico tratante. \\
Objetivo de tarea & Completar teleconsulta y recibir indicaciones clínicas. \\
Contexto & Domicilio, dispositivo móvil o computador y conectividad variable. \\
Evidencia & Eventos de sesión y resultado de finalización, con semilla 25022. \\
Característica principal & Cobertura del contexto; impactos secundarios en efectividad y satisfacción. \\
Métrica & M-CC-02: porcentaje de teleconsultas no completadas. \\
Resultado & 8,68\% de caídas frente a una meta inferior al 5\%. \\
Riesgo & Pérdida de continuidad, repetición de atención y menor confianza. \\
Recomendación & Reconexión, canal alterno y medición segmentada por red y dispositivo. \\
\bottomrule
\end{tabularx}
\end{table}

\subsection{Especificación de la medida}

Sea $S$ el conjunto de teleconsultas iniciadas y $C$ el subconjunto que termina correctamente. La tasa de caída se calcula como:

\begin{equation}
M\text{-}CC\text{-}02 = \frac{|S|-|C|}{|S|}\times 100.
\end{equation}

El resultado es verde si no supera 5\%, amarillo si está entre 5\% y 10\%, y rojo por encima de 10\%. Con 8,68\%, la línea base queda en amarillo. Esta clasificación expresa una alerta, no una aceptación clínica definitiva.

\subsection{Escenarios de prueba}

\begin{table}[H]
\centering
\caption{Cobertura propuesta para Telemedicina 2.0}
\begin{tabularx}{\textwidth}{p{2.6cm}p{2.8cm}YY}
\toprule
Contexto & Condición & Resultado esperado & Evidencia \\
\midrule
Red estable & WiFi con latencia baja & Sesión completa sin reconexión & Duración y finalización \\
Red degradada & Pérdida breve de paquetes & Recuperación sin perder la consulta & Intentos de reconexión \\
Cambio de red & WiFi a datos móviles & Continuidad de sesión o reingreso guiado & Cambio de contexto \\
Interrupción larga & Sin conectividad temporal & Canal alterno y conservación del estado & Motivo y recuperación \\
Dispositivo limitado & Cámara o audio no disponible & Diagnóstico previo y alternativa & Resultado de prechequeo \\
\bottomrule
\end{tabularx}
\end{table}

La recomendación no se aplica en la simulación inicial. Primero se conserva el 8,68\% como referencia. En una versión posterior se implementarían recuperación de sesión, prueba previa y canal alterno; después se repetiría la misma medida con segmentos por dispositivo, sede asociada y condición de red.

\subsection{Criterio de aceptación del siguiente ciclo}

La mejora se considerará válida si M-CC-02 queda por debajo del 5\%, no aumenta la duración hasta un nivel inaceptable, mantiene la privacidad de la sesión y mejora o conserva el CSAT. Así se evita resolver una dimensión perjudicando eficiencia, satisfacción o libertad de riesgo.
